\bigskip
\section{Introduction}\label{sec:intro}

\subsection{The problem: the model is asked to do three jobs it does poorly}

People increasingly ask general-purpose AI agents about their own sleep, heart rate and training. Today an agent
reaches wearable data in one of two ways. A connector exports the record and the model reads it in its context,
or a tool server, often a Model Context Protocol (MCP) server, lets the model fetch whatever its credential
reaches. One design puts the full record in front of the model; the other puts all of it one call away. Both leave
three jobs to the model: computing over the numbers, deciding what to disclose, and staying consistent from one
question to the next.

Each of these jobs is fragile when a language model does it. Language models reason poorly over raw numeric
series~\citep{merrill2024timeseries}, and their errors on compositional arithmetic compound with the number of
steps~\citep{dziri2023faith}; the Personal Health Insights Agent (PHIA) study found that an agent which writes and
runs code answers objective questions over wearable data far better than direct numerical reasoning
does~\citep{merrill2024phia}. Agents disclose information that is irrelevant to the task they were
given~\citep{shao2024privacylens,wang2025privacyinaction}, and text injected into the data they read can steer
them~\citep{greshake2023indirect}. Repeated questions over the same data yield different
answers~\citep{atil2024nondeterminism}.

Before any of these failures there is an earlier one. Privacy evaluations of language-model agents mostly ask
whether the model \emph{leaks}: whether sensitive information in its context reaches an output where it should
not~\citep{mireshghallah2024confaide,shao2024privacylens}. We argue that the step before is already a harm. The
principle at stake is \emph{data minimization}. In the United States, the ``minimum necessary'' standard of the
Health Insurance Portability and Accountability Act (HIPAA) binds covered entities and their business associates to
limit protected health information to what the intended purpose requires, with disclosures to the individual
exempt~\citep{hipaa164502}. We borrow the principle, although consumer wearable data mostly falls outside HIPAA,
under the Federal Trade Commission (FTC) Health Breach Notification Rule, which treats an unauthorized disclosure by a
health app as a breach~\citep{ftc2024hbnr}, and state consumer-health-data laws such as Washington's My Health My
Data Act, which requires consent to collect or share such data~\citep{wa2023mhmda}. API security has long treated an interface that returns more fields than its client needs, and relies on the client
to filter them, as a vulnerability in its own right: Excessive Data Exposure~\citep{owasp2019api3}, since folded into
Broken Object Property Level Authorization~\citep{owasp2023api3}. An agent given a raw wearable record is exactly
such a client. By analogy we name the failure \emph{context over-exposure}: the model's context receives personal
data beyond what the question requires. A question about last night's sleep does not need six weeks of heart rate,
workouts and body weight, yet a connector delivers all of it and a tool server makes all of it one call away.

We keep \emph{exposure}, a property of the model's input, distinct from \emph{leakage}, a property of its output.
Exposure matters because models are unreliable filters. Two strong models leaked in 25.68\% to 38.69\% of cases even
when instructed to protect privacy~\citep{shao2024privacylens}; contextual-integrity violations rise as more
entries enter the context~\citep{fu2026ciwork} and accumulate with use~\citep{mireshghallah2025cimemories}; and
anything in the context can be extracted by injected instructions~\citep{greshake2023indirect,bagdasarian2024airgap}.
Once delivered, the data also sits in the model provider's context, outside the user's control, whatever the model
does next. \sys{} therefore minimizes what enters the context, not only what leaves it.

\subsection{Approach}

We present \sys{}, a broker that sits between wearable data and any agent and takes the three jobs away from the
model (\cref{fig:arch}). It has three parts:
\begin{itemize}[leftmargin=1.4em,itemsep=2pt]
\item \textbf{Per-question retrieval by computation.} A large language model (LLM) planner turns each question into
  a typed query, a \emph{QuerySpec}, over a metric layer. Fixed, tested code computes the answer. The planner never
  sees the stored record, only the question text, and the model that reads the answer never computes it.
\item \textbf{Deterministic minimum-necessary release.} A policy gate outside the model checks every query against
  a grant that is specific to one agent and one purpose, and releases the coarsest of four levels (flag, bucket, values, raw) that answers it. It writes a hash-chained
  receipt for every release and every denial, and can cap the number of packets an agent receives per day.
\item \textbf{One response contract.} Every agent receives the same \emph{packet}: findings relative to the
  user's own baseline, data-quality caveats, freshness and sources, in at most 400 tokens.
\end{itemize}

We set out to measure what this design changes when the agent model is held fixed. We compare three conditions
that differ only in how the agent reaches the data: \condA{}, where the user's export is in the prompt; \condB{},
where unrestricted data tools return whatever is asked; and \condC{}, where one tool sends the question through the
broker. The comparison covers \nModels{} models, from open-weight models with a few billion active parameters to
frontier systems, on a synthetic cohort whose ground truth is recorded by construction and which contains the data-quality
traps real wearable data contains: duplicate step sources, time-zone travel, daylight-saving changes, missing
days, stale syncs and device switches. We measure accuracy, honest refusal, tokens, latency and cost, leakage under
contextual-integrity and injection scenarios, and repeatability. We repeat the comparison end to end on PHIA's
objective questions and on real Fitbit data from two public datasets, and evaluate the planner alone and the compute
engine alone on PHIA.

\subsection{Contributions}

\begin{itemize}[leftmargin=1.4em,itemsep=2pt]
\item \sys{}, a broker that separates planning (an LLM that sees only the question and a metric catalog),
  computation (tested code over a metric layer that handles duplicate sources, local days, coverage and personal
  baselines) and release (a deterministic gate with purpose-bound grants, four release levels and hash-chained
  receipts) (\cref{sec:design}).
\item A controlled comparison of \sys{} with the two prevailing designs on \nModels{} agent models, scored against
  ground truth rather than against the broker's own output, with paired exact tests and a bootstrap that resamples
  people (\cref{sec:study,sec:results}).
\item A measurement of context over-exposure, the share of the record each design puts in front of the model per
  question, kept separate from leakage, which we score with a detector checked against hand labels
  (\cref{sec:res:exposure,sec:res:privacy}).
\item An account of where the benefit lies and where \sys{} fails, including end-to-end results on PHIA and real
  Fitbit data, and two pre-registered changes that address the design decisions PHIA exposed
  (\cref{sec:res:accuracy,sec:res:failures,sec:res:bench,sec:res:v2}).
\end{itemize}

Computing outside the model~\citep{gao2022pal,chen2023pot,merrill2024phia} and deterministic policy outside the
model~\citep{shi2025progent,debenedetti2025camel,kodathala2026aiauthz} each exist; our contribution is their
combination for personal wearable data, with release levels and purpose-bound grants, and the measurement of what it
changes across models.
